\newpage
\begin{center}
\large\textbf{Минобрнауки России}

\large\textbf{Юго-Западный государственный университет}
\vskip 1em
\normalsize{Кафедра программной инженерии}
\vskip 1em
\ifВКР{
        \begin{flushright}
        \begin{tabular}{p{.4\textwidth}}
        \centrow УТВЕРЖДАЮ: \\
        \centrow Заведующий кафедрой \\
        \hrulefill \\
        \setarstrut{\footnotesize}
        \centrow\footnotesize{(подпись, инициалы, фамилия)}\\
        \restorearstrut
        «\underline{\hspace{1cm}}»
        \underline{\hspace{3cm}}
        20\underline{\hspace{1cm}} г.\\
        \end{tabular}
        \end{flushright}
        }\fi
\end{center}
\vspace{1em}
  \begin{center}
  \large
\ifВКР{
ЗАДАНИЕ НА ВЫПУСКНУЮ КВАЛИФИКАЦИОННУЮ РАБОТУ
  ПО ПРОГРАММЕ БАКАЛАВРИАТА}
  \else
ЗАДАНИЕ НА КУРСОВУЮ РАБОТУ (ПРОЕКТ)
\fi
\normalsize
  \end{center}
\vspace{1em}
{\parindent0pt
  Студента \АвторРод, шифр\ \Шифр, группа \Группа
  
1. Тема «\Тема\ \ТемаВтораяСтрока»
\ifВКР{
утверждена приказом ректора ЮЗГУ от \ДатаПриказа\ № \НомерПриказа
}\fi.

2. Срок предоставления работы к защите \СрокПредоставления

3. Исходные данные для создания программной системы:

3.1. Перечень решаемых задач:

\begin{enumerate}
	\item Спроектировать архитектуру серверной части корпоративного мессенджера с использованием Python, WebOb и WSGI, обеспечив маршрутизацию запросов и модульную структуру приложения.
	
	\item Реализовать механизм регистрации, авторизации и управления пользовательскими сессиями с привязкой к уникальному идентификатору и криптографическим ключам, обеспечив безопасность передачи сообщений.
	
	\item Разработать интерфейсы (API) для обработки запросов, связанных с отправкой, получением, редактированием и удалением сообщений в групповых и приватных чатах, включая асинхронную проверку обновлений.
	
	\item Организовать работу с файловыми вложениями: загрузка, сохранение, определение MIME-типов, отображение изображений и предоставление возможности скачивания других типов файлов.
	
	\item Внедрить модуль управления группами: создание, переименование, добавление и исключение участников, управление ролями (участник, администратор, владелец), проверка прав доступа и уведомления о событиях.
	
	\item Обеспечить серверную обработку системных сообщений, их отображение в интерфейсе.
	
	\item Разработать механизм полнотекстового поиска по сообщениям с поддержкой фильтрации по типу чата, ID чата и сортировкой по дате или релевантности.
	
	\item Спроектировать обработку ошибок с возвратом корректных кодов HTTP-ответов, логированием исключений и сохранением отказоустойчивости сервиса.
	
	\item Провести модульное и интеграционное тестирование логики серверной части: проверка корректности маршрутов, безопасности доступа, устойчивости к некорректным запросам и проверка взаимодействия с базой данных.
\end{enumerate}



{\parindent0pt
  3.2. Входные данные и требуемые результаты для программы:}

\begin{enumerate}
	\item Входными данными для серверной части корпоративного мессенджера являются: HTTP-запросы от клиентской части (регистрация, авторизация, отправка сообщений, действия над группами и вложениями); данные о пользователях, сессиях, сообщениях и ролях, получаемые из базы данных; файлы, прикреплённые к сообщениям (изображения, документы и прочее); параметры поиска и фильтрации сообщений; данные cookie и параметров сессии для идентификации пользователей.
	
	\item Выходными данными являются: сформированные ответы на упорядоченные коллекции сообщений (групповых, приватных, общего чата); данные о пользователях (идентификаторы, роли, активность); структуру групп и участников; сведения о вложениях (пути к файлам, MIME-типы, имена); статусы операций; временные метки и параметры синхронизации.
\end{enumerate}


{\parindent0pt

  4. Содержание работы (по разделам):
  
  4.1. Введение.
  
  4.1. Анализ предметной области.
  
4.2. Техническое задание: основание для разработки, назначение разработки,
требования к программной системе, требования к оформлению документации.

4.3. Технический проект: общие сведения о программной системе, проект
данных программной системы, проектирование архитектуры программной системы, проектирование пользовательского интерфейса программной системы.

4.4. Рабочий проект: спецификация компонентов и классов программной системы, тестирование программной системы, сборка компонентов программной системы.

4.5. Заключение.

4.6. Список использованных источников.

5. Перечень графического материала:

\списокПлакатов

\vskip 2em
\begin{tabular}{p{6.8cm}C{3.8cm}C{4.8cm}}
Руководитель \ifВКР{ВКР}\else работы (проекта) \fi & \lhrulefill{\fill} & \fillcenter\Руководитель\\
\setarstrut{\footnotesize}
& \footnotesize{(подпись, дата)} & \footnotesize{(инициалы, фамилия)}\\
\restorearstrut
Задание принял к исполнению & \lhrulefill{\fill} & \fillcenter\Автор\\
\setarstrut{\footnotesize}
& \footnotesize{(подпись, дата)} & \footnotesize{(инициалы, фамилия)}\\
\restorearstrut
\end{tabular}
}

\renewcommand\labelenumi{\theenumi.}
